\section{Conclusions and Future Work}

In this paper we identified a previously undescribed mechanism for relational composition in transformers, which we call relative magnitude-based composition: the model encodes bigram order through the relative sizes of scalar attention weights rather than through distinct directions in the activation space. 
We characterised this mechanism mathematically via circuit  analysis, showed that SAEs trained on this model reliably learn graded special latents whose activation magnitudes correlate with $\Delta\alpha$, and demonstrated that steering these latents reverses the model's predicted output for over 50\% of inputs. 
These results indicate that the activation values of SAE latents, and not only their active/inactive status, carry information that is necessary for understanding model behaviour in this setting, which is a different perspective on transformer features than is currently prevalent in the mechanistic interpretability literature.
We release our trained transformers and SAEs to support future work and reproducibility of our results here: \url{https://wandb.ai/theo-farrell99-durham-university/order-by-scale}, \url{https://wandb.ai/theo-farrell99-durham-university/orderbyscale_sae_sweep}.


\subsection*{Limitations}

Our mechanism analysis is conducted on a single minimal toy model with a narrow, controlled task. The architecture omits components standard in pretrained LLMs (MLPs, LayerNorm, multi-head attention, value and output projections) and uses a vocabulary of only 100 symbols. However, our SAE experiments demonstrate robust findings across diverse configurations: special latents (characterized by $\Delta\alpha$ correlation) emerge across three distinct SAE architectures (BatchTopK, JumpReLU, Matryoshka) with varying sparsity and dictionary size settings \ref{tab:sae_sweep_configs}. This suggests the mechanism may be fundamental to how SAEs represent the relational composition in this model, though whether relative magnitude-based composition appears in less constrained models remains an open question.

A further limitation is that our circuit analysis characterises the mechanism on inputs the model solves correctly. The roughly 9\% of inputs on which the model errs seem to reflect cases where $\Delta\alpha \approx 0$ (as shown in Figure~\ref{fig:confusion-band}), where the magnitude-based mechanism cannot resolve the two orderings. A complete account of the model's behaviour would need to explain these failure cases as well.

Additionally, the steering pipeline currently handles one special latent at a time; simultaneous multi-latent steering could cause even more outputs to be swapped, as suggested by the co-activation experiments in Section~\ref{ss:steering_exp}. Finally, our SAE checkpoint was selected on the basis of having few dead latents, which introduces selection bias toward high-quality SAEs. While generalisation experiments (Section~\ref{ss:sae_generalisation}) show key findings hold broadly, a fully representative study would require analysis across random samples of all trained SAEs.

\subsection*{Future Work}

The most important next step is determining whether pretrained LLMs implement similar mechanisms. The GemmaScope suites of SAEs \citep{lieberum2024gemma, mcdougall2025gemmascope2} provide a natural test environment, as it includes JumpReLU and Matryoshka SAEs trained on Gemma 2 and Gemma 3 at multiple widths - the same SAE variants we study here. The specific target would be identifying SAE latents in pretrained models where activation magnitude, rather than feature presence, is causally relevant for downstream behaviour.

There are also several remaining extensions within the toy model setting. Scaling to larger $N$-grams would test how the minimum number of layers required for high accuracy scales with $N$, and whether the magnitude-based mechanism persists or is replaced by a direction-based alternative as the task becomes harder. 
Moreover, the role of positional encodings in supporting this mechanism is not fully understood, and it is unclear whether relative magnitude-based composition is possible under alternative positional encoding schemes such as absolute encodings \citep{radford2019language} or RoPE \citep{su2024roformer}. 
% In our study we sampled $d_i$ uniformly from $[0,99]$. In future, not only will we extend our research to larger ranges, but also analyse how or whether our results change if $d_i$ are distributed normally in those ranges rather than uniformly. We may further extend this to other distributions if it seems fruitful.


Furthermore, it is an open question as to whether the pathfinding models trained by \citet{brinkmann2024mechanistic}, which were the inspiration for this paper, learn a similar relation composition mechanism.

On the SAE side, the steering pipeline could be extended in two 
directions: improving computational efficiency to allow exhaustive testing across all trained SAE configurations rather than a cherry-picked sample, and testing whether scaling multiple latents simultaneously can recover the remaining inputs for which single-latent steering fails. 
Preliminary evidence in Section~\ref{ss:steering_exp} suggests this may be possible. 
Finally, the failure modes of the base transformer model (roughly 9\% of input) have not been fully characterised, and it is worth investigating whether these relate systematically to the cases where latent steering also fails.

In summary, the findings in this paper are intended as a foundation for further research on magnitude-based relational composition mechanisms in LLMs.